\section{引言}\srcpage{1}

Langlands 纲领是一张影响深远的猜想与定理之网，它在数学的许多不同领域中都有化身。虽然它最初由 Langlands 提出，用以处理数论中的问题，例如推广类域论的非交换互反律的存在性，以及算术局部对称空间上的自守形式理论。但此后它又在许多不同的情形中找到了化身。这些翻译通常依照 Weil 的「Rosetta Stone」（罗塞塔石碑）进行，它有三栏：
\begin{enumerate}
\item 数域的原初情形；
\item 有限域 $\Fq$ 上函数域的、与之类似而更易接近的情形；等价地，$\Fq$ 上光滑射影曲线 $X$ 的情形；
\item 代数闭域上光滑射影曲线的几何情形；例如在 $\C$ 上，这些等价于 Riemann 曲面。
\end{enumerate}

\srcpage{2}
从 1970 年代 Drinfeld 的工作开始，人们研究了 $\Fq$ 上以及其代数闭包 $\overline{\F}_q$ 上光滑射影曲线 $X$ 的 Langlands 纲领。此后数年间，人们也研究了 $\C$ 上光滑射影曲线的类似结构；但有一段时间，这些研究似乎离最初关心的数论情形越来越远，同时又发现了与数学物理的新联系。当我还是研究生、在算术 Langlands 纲领中工作时，我无法理解几何 Langlands 纲领所研究的结构究竟是什么，为什么人们应当对它们感兴趣，以及它们与最初关心的问题有何关联。

然而，在过去几年里，情况发生了剧变：（特征 $0$、处处非分歧、整体）几何 Langlands 猜想已被证明。并且，这一几何猜想又被联系回 Weil 罗塞塔石碑的第二栏，即有限域上的函数域，从而对函数域上 Langlands 纲领中一些最基本的问题给出非常精确的回答。

在本引言的余下部分，我将简要回顾 Langlands 最初在数域情形提出的问题，然后陈述 Gaitsgory–Raskin（2025a）证明的一个定理。该定理依赖许多其他论文，作者包括 Dima Arinkin、Dario Beraldo、Justin Campbell、Lin Chen、Joakim Færgeman、Dennis Gaitsgory、David Kazhdan、Kevin Lin、Sam Raskin、Nick Rozenblyum、Yakov Varshavsky——它对该函数域情形中的类似问题给出了非常精确的回答。

就其核心而言，Langlands 纲领试图描述自守函数构成的向量空间。这不过是对某个极为基本之物给出的一个花哨名称。也就是说，数学中出现的一些最基本、最突出、最漂亮的空间例子，正是那些局部以某个给定几何（例如双曲几何或球面几何）为模型的空间。更确切地说，人们固定一个 Riemann 对称空间 $M$\footnote{习惯上把该对称空间记为 $X$，但这与几何 Langlands 中所用曲线的标准记号冲突。}及其自同构李群 $G$，然后考虑局部同构于 $M$ 的流形 $N$，其坐标变换属于 $G$。若 $N$ 连通且完备，则它可以写成局部对称空间 $M/\Gamma$，其中 $\Gamma=\pi_1(N)$ 通过一个嵌入作为 $G$ 的离散子群作用。著名的刚性定理表明，这一情形实际上在大多数情况下比它看起来更具代数与数论色彩——也就是说，$G$ 带有约化线性代数群的结构，而 $\Gamma\subset G$ 是算术群。

除了一些细节，自守函数就是算术群 $\Gamma$ 对应的局部对称空间 $M/\Gamma$ 上的函数。在本引言中，我将忽略所考虑的函数空间究竟为何这一问题，而直接把 $M/\Gamma$ 上某个合适的性质良好的函数构成的向量空间记为 $\cA(M/\Gamma)$。注意，李群 $G$ 作用在 $M$ 上；无穷小作用特别地通过微分给出李代数 $\gc$ 在 $\cA(M)$ 上的作用。然而，这一作用不与 $\Gamma$ 交换，因此不在 $\cA(M/\Gamma)$ 上诱导作用。但我们可以把 $\gc$-作用延拓到泛包络代数 $U(\gc)$ 上。由

\srcpage{3}
Harish-Chandra 的工作，$U(\gc)$ 有很大的中心 $Z(U(\gc))$，它同构于 Cartan 子代数 $\mathfrak{h}\subset\gc$ 的泛包络代数中的 Weyl 群不变量 $U(\mathfrak{h})^W$。由于它的中心性，$Z(U(\gc))$ 在 $\cA(M)$ 上的作用现在与 $G$ 的作用、从而与 $\Gamma$ 的作用交换，于是它在 $\cA(M/\Gamma)$ 上诱导作用。特别地，这给出 Laplace 算子的作用。

另一方面，由于 $\Gamma$ 是算术的，它与许多共轭子群可公度。为简化讨论，让我作如下假设。

\begin{assumption}
群 $G$ 分裂且单连通。从现在起，我将用 $G$ 记它在 $\Z$ 上的分裂形式，于是先前的 $G$ 现在称为 $G(\R)$。此外，$\Gamma=G(\Z)$ 是 $G$ 的全部整点构成的群。
\end{assumption}

对任意 $g\in G(\Q)$，群 $\Gamma_g=\Gamma\cap g^{-1}\Gamma g$ 是 $\Gamma$ 的有限指数子群。此时，$g$ 在 $M$ 上的作用下降为一个映射
\[ M/\Gamma_g\xrightarrow{g}M/\Gamma, \]
连同自然投影 $M/\Gamma_g\to M/\Gamma$，我们得到 Hecke 对应
\[ \begin{tikzcd} & M/\Gamma_g \arrow[dl] \arrow[dr, "g"] & \\ M/\Gamma & & M/\Gamma. \end{tikzcd} \]
这诱导出自同态
\[ T_g\colon\cA(M/\Gamma)\to\cA(M/\Gamma), \]
其做法是先把函数拉回到 $M/\Gamma_g$，再沿有限映射 $g\colon M/\Gamma_g\to M/\Gamma$ 的纤维求和。这些 Hecke 算子与上面构造的微分算子交换，而且有些出人意料的是，它们彼此也交换\footnote{这里用到 $\Gamma=G(\Z)$；一般情形略为复杂。}。合起来，它们给出一个（无限生成的）交换 $\C$-代数 $\mathbb{T}=\mathbb{T}_G$ 在 $\cA(M/\Gamma)$ 上的作用。

下述问题可以说是 Langlands 纲领的中心问题。

\begin{problem}
描述带有 $\mathbb{T}$ 作用的向量空间 $\cA(M/\Gamma)$。
\end{problem}

Langlands 的进路是把 $\cA(M/\Gamma)$ 作为 $\mathbb{T}$-模作谱分解。他证明它可分解为「连续谱」与「离散谱」。此外，连续谱允许用所谓的 Eisenstein 级数给出一个（复杂但本质上是显式的）描述，这些级数由与 $G$ 的真 Levi 子群相伴的局部对称空间上的自守函数归纳地构造。「离散谱」允许分解为不可约「自守表示」的直和。其中一些仍间接地 \srcpage{4}来自 Eisenstein 级数，作为「Eisenstein 级数的留数」，而另一些则是「对 $G$ 真正新的」；这些被称为「尖点的」。

关于可能的「Hecke 特征值系统」，即 $\cA(M/\Gamma)$ 的支撑中的同态 $\psi\colon\mathbb{T}\to\C$，存在若干深刻的预测。为了陈述它们，多谈一点 Hecke 算子的结构会有用。

我们关心空间 $M/G(\Z)$，其中 $M=G(\R)/K_\infty$ 是 $G(\R)$ 的对称空间；这里 $K_\infty\subset G(\R)$ 是极大紧子群。给出该商的一个点，等价于给出一个 $G(\Z)$-挠子——等价地，给出 $\Spec(\Z)$ 上的一个 $G$-挠子 $E$——连同相伴的 $G(\R)$-挠子到 $K_\infty$ 的一个约化——换言之，$E\times_{\Spec(\Z)}\Spec(\R)$ 上的一个「度量」。例如，若 $G=\mathrm{SL}_n$，从而 $K_\infty=\mathrm{SO}_n(\R)\subset\mathrm{SL}_n(\R)$，则这是如下数据：秩为 $n$ 的 $\Z$-格 $E$，连同平凡化的行列式 $\det(E)\cong\Z$，以及 $E\otimes_\Z\R$ 上的一个正定配对。

在这个图像中，由 $g\in G(\Q)$ 标号的 Hecke 算子一般地不改变 $G$-挠子 $E$，即在 $\Spec(\Q)$ 上不改变它，而是在 $\Spec(\Z)$ 的某些有限位处修改它。实际上，这可以一次只对一个素数 $p$ 进行，使用 $g\in G(\Z[\frac1p])$ 的 Hecke 算子。$p$ 处的 Hecke 算子给出代数
\[ \C\big[G(\Z)\backslash G(\Z[\tfrac1p])/G(\Z)\big] \]
的作用，该代数由 $G(\Z[\frac1p])$ 中 $G(\Z)$-双陪集上的紧支撑函数组成，其中 $T_g$ 对应双陪集 $G(\Z)gG(\Z)$ 的特征函数。实际上，由强逼近，该代数可以用 $p$ 进群 $G(\Q_p)$ 改写成
\[ \C\big[G(\Z_p)\backslash G(\Q_p)/G(\Z_p)\big]. \]

Harish-Chandra 对泛包络代数中心的描述，其类似物是 Satake 同构：固定分裂极大环面 $T\subset G$ 及其 Weyl 群 $W$，则有 $\C$-代数同构
\[ \C\big[G(\Z_p)\backslash G(\Q_p)/G(\Z_p)\big]\cong\C[X_*(T)]^W; \]
特别地，左端交换。人们常常形式地引入 $\C$ 上的对偶环面 $\Td$，其特征群 $X^*(\Td)=X_*(T)$ 与 $T$ 的余特征群等同；于是
\[ \C[X_*(T)]^W=\cO(\Td/W), \]
即仿射簇 $\Td/W$ 上正则函数构成的代数。若再引入 Langlands 对偶群 $\Gd$，其根数据与 $G$ 的根数据对偶，特别地它以 $\Td$ 为极大环面、以 $\check W=W$ 为 Weyl 群，则我们可以进一步把它重新解释为
\[ \cO(\Gd/_{\mathrm{ad}}\Gd), \]
即 $\Gd$ 上的共轭不变函数。因此，Satake 同构给出同构
\[ \C\big[G(\Z_p)\backslash G(\Q_p)/G(\Z_p)\big]\cong\cO(\Gd/_{\mathrm{ad}}\Gd), \]
这是群 $G$ 与其 Langlands 对偶群 $\Gd$ 之间在函数层面的第一座桥梁。

\srcpage{5}
上述讨论给出抽象 Hecke 代数
\[ \mathbb{T}=\bigotimes_{p}\C\big[G(\Z_p)\backslash G(\Q_p)/G(\Z_p)\big]\otimes Z(U(\gc)) \]
在 $\cA(M/\Gamma)$ 上的作用。考虑 $\cA(M/\Gamma)$ 的支撑中的一个映射 $\psi\colon\mathbb{T}\to\C$。映射 $\psi$ 对应于对所有 $p$ 的映射
\[ \C\big[G(\Z_p)\backslash G(\Q_p)/G(\Z_p)\big]\to\C, \]
以及映射 $Z(U(\gc))\to\C$。由 Satake 同构与 Harish-Chandra 同构，它们分别给出：对每个素数 $p$，一个 $\Gd(\C)$ 中元素的 $\Gd$-共轭类 $\psi_p$；以及 $\psi_\infty\in\check{\gc}$。

有两个核心猜想；通过归纳，可以假设 $\psi$ 属于离散谱。在陈述下述猜想时，我先假设 $G=\mathrm{SL}_n$（从而 $\Gd=\mathrm{PGL}_n$），且 $\psi$ 是尖点的。
\begin{enumerate}
\item （Ramanujan）对每个 $p<\infty$，$\psi_p$ 落在 $\Gd(\C)$ 的一个紧子群中；而 $\psi_\infty$ 落在 $\Gd(\C)$ 的某个紧子群的李代数中。
\item （互反）$\Spec(\Z)$ 上存在不可约 $\Gd$-局部系统\footnote{警告：这一概念尚不知道，参看下面的注 1.4。}，其中 $\psi_p$ 是沿「圆」$\Spec(\Fp)\subset\Spec(\Z)$ 的单值性，而 $\psi_\infty$ 给出广义 Hodge 数。
\end{enumerate}
对完整的离散谱，以及对 $\mathrm{SL}_n$ 以外的群，这两个猜想都需要稍作修改；Arthur（1989）陈述了相应的推广。

\begin{remark}
Ramanujan 猜想是一个极其简单而一般的陈述。应用于 $G=\mathrm{SL}_2$，并且只看微分算子的作用，它预测在 $\mathbb{H}/\mathrm{SL}_2(\Z)$ 上 Laplace 算子的谱隙至少为 $\frac14$；实际上，它对所有算术双曲曲面都预测同样的事。此外，它把这一现象适当地推广到所有算术局部对称空间。对 $p<\infty$ 的 $\psi_p$，它预测 Hecke 特征值的「平方根消去」；在某些情形，Deligne 通过证明一个互反的版本与 Weil 猜想，即「有限域上的 Riemann 假设」，著名地证明了这一点——特别是在 Ramanujan 最初提出其猜想的权 $12$ 模形式 $\Delta$ 的情形。

注意，Ramanujan 猜想非常强烈地提示存在某个与 $\psi$ 相伴的 $\Gd(\C)$ 的典范紧子群，在共轭意义下所有 $\psi_p$ 都落在其中（而 $\psi_\infty$ 落在其李代数中）。实际上，人们预期这样的子群存在，且 $\psi_p$ 在其中等分布——这就是 Sato–Tate 猜想。此外，这个紧子群应当恰好是互反猜想所预测的 $\Gd$-局部系统的单值群。
\end{remark}

\begin{remark}
遗憾的是，互反猜想就没有同样的简单性：它只是一个元猜想。也就是说，相关的 $\Gd$-局部系统概念尚不知道。定义这样一个概念，本质上等价于定义 $\Spec(\Z)$ 上 $\C$-局部系统的某个 Tannakian 范畴，而这又等价于定义

\srcpage{6}
某个投射代数（或投射紧）群——这就是猜想中的 Langlands 群。然而，我相信人们不应当把注意力集中在 Langlands 群上，而应当寻找 $\Spec(\Z)$ 上 $\C$-局部系统的一个更一般的理论——或者说，真正要对 $\Spec(\Z)$ 上的概形建立「扭量结构的变体」的完整六函子形式化。这样的理论应当允许 Hodge 数是复数而非整数——在 $\Spec(\R)$ 上，这在 Scholze（2024）发展的理论中是可行的。

若 $\psi_\infty\in\check{\gc}$ 的共轭类在 $\check{\tc}$ 中有整的代表，从而所预测的 Hodge 数是整数，则可以使用一个非猜想的替代物：也就是说，固定某个素数 $\ell$ 与同构 $\C\cong\overline{\Q}_\ell$，并使用 étale $\overline{\Q}_\ell$-局部系统。回忆一下，$\ell$ 进局部系统本质上是 $\ell$ 进 Galois 表示，这就给出自守形式与 Galois 表示之间著名的猜想性关系\footnote{此处我忽略了一个重要的 $\rho$-shift（$\rho$-平移），关于这一点的详细阐述参看 Buzzard–Gee（2014）。}。
\end{remark}

上述猜想有一系列加强版本，它们不仅预测某个 Hecke 特征值系统 $\psi$ 的支撑，还预测它出现的重数。特别地，这预测在某种意义上 $\Spec(\Z)$ 上所有 $\Gd$-局部系统都应当来自某个自守形式。

下面的猜想实际上会对问题 1.2 给出更精确的回答。如同 Langlands 纲领中许多最深刻的猜想一样，它其实是一个猜想的模板——许多术语仍未定义\footnote{此外，还缺少一些精确化；尤其是，阿基米德位处的处理可能需要一些修正，而且我忽略了由各种 $\rho$-shift 引起的可能扭转。}。

\begin{conjecture}
存在 $\C$ 上的一个叠 $\LocSys_{\Gd}$，它参数化 $\Spec(\Z)$ 上的 $\Gd$-局部系统，并带有一个映射 $\LocSys_{\Gd}\to\Spec(\mathbb{T})$，该映射由把 $\Gd$-局部系统限制到所有 $\Spec(\Fp)\subset\Spec(\Z)$、以及广义 Hodge 数诱导。此外，存在自守函数空间与 $\LocSys_{\Gd}$ 上对偶复形 $\omega_{\LocSys_{\Gd}}$ 的整体截面之间的典范 $\mathbb{T}$-等变同构
\[ \cA(M/\Gamma)\cong\Gamma(\LocSys_{\Gd},\omega_{\LocSys_{\Gd}}). \]
\end{conjecture}

把这个猜想拆成两部分是有益的：
\begin{enumerate}
\item 可以把 $\cA(M/\Gamma)$ 提升为 $\LocSys_{\Gd}$ 上的拟凝聚层。
\item 这个拟凝聚层就是对偶复形。
\end{enumerate}
第一部分已经给出 Langlands 互反的「自守到 Galois」方向，它蕴含 $\cA(M/\Gamma)$ 的支撑必须包含在 $\LocSys_{\Gd}\to\Spec(\mathbb{T})$ 的像中。第二部分则相当于描述给定的 $\Gd$-局部系统出现的重数，并包含「Galois 到自守」的反方向。

\begin{remark}
该猜想的以下方面尚不清楚：

\srcpage{7}
\begin{enumerate}
\item $\Gd$-局部系统的概念是什么？
\item $\LocSys_{\Gd}$ 是哪种「叠」？
\end{enumerate}
我们注意到，这里肯定需要某种解析几何，这已经是因为：所寻求的向量空间 $\cA(M/\Gamma)$ 带有非平凡拓扑，因此必须从泛函分析的观点来研究。
\end{remark}

受下文将讨论的几何 Langlands 联系的启发，近年来若干作者群体提出了类似的猜想，例如 Zhu（2025，第 4.7 节）与 Emerton。不过，通常不是针对通常的自守形式，而是针对 $p$ 进自守形式（例如以 $M/\Gamma$ 的 $p$ 进奇异上同调的形式出现）来表述。在那种情形，所需的 $\Gd$-局部系统模空间可以（本质上）\footnote{即便在那里，阿基米德位也可能引起一些问题；不过，可以用 Clausen 最近的 Weil–Moore anima 理论来修正。}用 $p$ 进 étale 局部系统来定义。

猜想 1.5 会比 Arthur–Langlands 的猜想更精确：也就是说，它用 $\Gd$-局部系统描述带有 $\mathbb{T}$-作用的实际的自守形式向量空间；而 Langlands 只着眼于描述出现的自守表示的同构类与重数。令人惊讶的是，Arthur 精细化的参数概念并不（也不应当）进入它的陈述。相反，当人们分析右端的离散谱究竟是什么时，明显的部分来自 $\LocSys_{\Gd}$ 的孤立点（这些对应不可约 $\Gd$-局部系统），但进一步的贡献来自对偶复形 $\omega$ 变得更复杂的奇点。对这一现象的分析应当能恢复 Arthur 对 Langlands 猜想的精细化\footnote{正如我将在下文讨论的，在函数域情形这会成为一个精确的问题，它在 Gaitsgory–V. Lafforgue–Raskin 最近的工作中得到处理，参看 Raskin（2025）。}。

Gaitsgory、Raskin 等人的工作处于这一情形的函数域类比之中。粗略地说，这相当于把 $\Z$ 换成 $\Fp[\mathbf{t}]$，把 $\R$ 换成 $\Fp((\mathbf{t}^{-1}))$。注意，正如 $\Z\subset\R$ 是离散余紧子群，$\Fp[\mathbf{t}]\subset\Fp((\mathbf{t}^{-1}))$ 也是离散余紧子群。$M=G(\R)/K_\infty$ 的类比是 $G(\Fp((\mathbf{t}^{-1})))/G(\Fp[[\mathbf{t}^{-1}]])$，其中 $G(\Fp((\mathbf{t}^{-1})))$ 是以 $G(\Fp[[\mathbf{t}^{-1}]])$ 为极大紧子群的局部紧群。$\Gamma=G(\Z)$ 的类比是 $G(\Fp[\mathbf{t}])$。我们于是被引向考察双商
\[ G(\Fp[\mathbf{t}])\backslash G(\Fp((\mathbf{t}^{-1})))/G(\Fp[[\mathbf{t}]]). \]
结果表明它可以有更简洁的描述：它就是 $\mathbb{P}^1_{\Fp}$ 上 $G$-挠子的同构类集合。事实上，这不过是在镜像我们先前用 $\Spec(\Z)$ 上的 $G$-挠子连同 $\Spec(\R)$ 上的一个「度量」对 $M/\Gamma$ 的描述。

正如 $\Z$ 上的 Langlands 纲领对任意数域 $K$ 有关于 $\cO_K$ 的变体，我们在函数域情形也可以允许 $\mathbb{P}^1_{\Fp}$ 的任意有限覆盖，即对 $\Fp$ 上任意光滑射影曲线 $X$ 工作。设 $\Fq$ 为其常数域，即 $\Fp$ 的有限扩张。由 Hecke 的一个定理的函数域类比（参看 Weil（1974，第 291 页）），人们总能找到「Theta 特征」，即 $X$ 上线丛 $\OmX$ 的一个平方根 $\Theta$，我们将其固定\footnote{这一点只在稍后的归一化中用到，参看注 1.8。}。

\srcpage{8}
自守函数空间的类比是函数空间 $\cA(\Bun_G(\Fq))$，即 $X$ 上 $G$-丛的模空间 $\Bun_G$ 的 $\Fq$-点上的函数空间；从第 2 节开始将更详细地讨论它。Langlands 纲领的几乎所有方面都忠实地翻译到函数域情形，但许多事情容易得多：例如，$\Bun_G(\Fq)$ 只是一个离散点集，因此作为抽象向量空间，$\cA(\Bun_G(\Fq))$ 的描述是平凡的——这已经去掉了所有非平凡的分析。但同样地，有 Hecke 算子作用在 $\cA(\Bun_G(\Fq))$ 上。更确切地说，对每个点 $x\in X$，我们得到 Hecke 对应
\[ \begin{tikzcd} & \Hck_{G,x}(\Fq) \arrow[dl, "p_1"'] \arrow[dr, "p_2"] & \\ \Bun_G(\Fq) & & \Bun_G(\Fq), \end{tikzcd} \]
其中 $\Hck_{G,x}$ 参数化 $G$-丛对 $E,E'$，连同它们在 $X\setminus\{x\}$ 上的一个同构。这在 $\cA(\Bun_G(\Fq))$ 上诱导出 Hecke 代数
\[ \mathbb{T}_x=\C\big[G(\cO_x)\backslash G(F_x)/G(\cO_x)\big]\cong\cO(\Gd/_{\mathrm{ad}}\Gd) \]
的作用。这里 $\cO_x$ 表示 $X$ 在 $x$ 处的完备局部环——一个完备离散赋值环——而 $F_x$ 是它的分式域；这完全类似于 $\Z_p\subset\Q_p$，其中 $x$ 扮演 $p$ 的角色。也就是说，$h\in\C[G(\cO_x)\backslash G(F_x)/G(\cO_x)]$ 通过如下公式作用：
\[ f\in\cA(\Bun_G(\Fq))\mapsto\Big(x\mapsto\sum_{y,\ p_2(y)=x}h(y)f(p_1(y))\Big)\in\cA(\Bun_G(\Fq)). \]

同样地，Arthur–Langlands 的猜想预测用 $X$ 上 $\Gd$-局部系统来描述。仍然没有 $\C$-局部系统的概念，但在函数域情形，人们可以直接使用 étale $\overline{\Q}_\ell$-局部系统的概念，任选一个同构 $\C\cong\overline{\Q}_\ell$——根本原因在于，在函数域情形，自守函数的概念是纯代数的。

Zhu（2025）的工作表明，存在一个性态良好的（导出 Artin）叠 $\LocSys_{\Gd}$，它定义在 $\overline{\Q}_\ell$ 上，参数化 $\Gd$-局部系统，并带有一个到 $\Spec(\mathbb{T})$ 的映射；这里 $\mathbb{T}=\bigotimes_{x\in X}\mathbb{T}_x$ 是抽象 Hecke 代数。下述定理（本质上）在函数域情形解决了猜想 1.5：

\begin{theorem}[Gaitsgory 与 Raskin，2025a，定理 0.2.2]
设 $G$ 是 $\Fq$ 上的约化群，并假设 $p$ 相对于 $G$ 充分大。则存在既开又闭的子叠 $\LocSys'_{\Gd}\subset\LocSys_{\Gd}$，以及典范的 $\mathbb{T}$-等变同构
\[ \cA_c(\Bun_G(\Fq))\cong\Gamma(\LocSys'_{\Gd},\omega_{\LocSys'_{\Gd}}). \]
\end{theorem}

\srcpage{9}
这里 $\cA_c$ 表示紧支撑函数。Raskin（2025）宣布了与 Gaitsgory、V. Lafforgue 的合作工作，在函数域情形推出 Arthur–Ramanujan 猜想。因此，归根结底，人们证明的是关于实际数字的具体陈述，即 Hecke 算子的特征值！

当然，人们预期可以取 $\LocSys'_{\Gd}=\LocSys_{\Gd}$；这对 $G=\mathrm{SL}_n$ 是已知的。

\begin{remark}[Whittaker 归一化]
这个定理有一个令人惊讶的方面：它给出向量空间之间一个特别的同构。归一化的一个关键方面是下述匹配，它用到我们对 Theta 特征 $\Theta$ 的选取，以及在 $\C$ 中单位根的标准选取。在右端，$\LocSys_{\Gd}$ 的几何性质（其余切复形的行列式典范平凡）保证 $\omega_{\LocSys_{\Gd}}$ 带有一个典范截面，于是右端有一个特别的向量。Gaitsgory–V. Lafforgue–Raskin 证明，这对应于 Poincaré 级数
\[ \mathcal{P}\in\cA_c(\Bun_G(\Fq)), \]
它是把 Whittaker 特征
\[ \Bun_U^{\Th}(\Fq)\to\C^\times \]
沿 $\Bun_U^{\Th}(\Fq)\to\Bun_G(\Fq)$ 的纤维求和而得到的。这里 $\Bun_U^{\Th}$ 是 $U$-丛模空间的一个扭转形式，其中 $U$ 是 Borel 子群 $B\subset G$ 的单幂根。更确切地说，$\Bun_U^{\Th}$ 是 $U^{\Th}$-丛的模空间，这里 $U^{\Th}$ 是 $X$ 上的一个群概形，即 $U$ 被 $T$-挠子 $2\rho(\Theta)$ 扭转所得的扭转形式，而该挠子是线丛 $\Theta$ 所对应的 $\Gm$-挠子在映射 $2\rho\colon\Gm\to T$ 下的前推。这一扭转保证：对任意单根 $\alpha$，$U^{\Th}$ 的相应根空间不是 $\mathbb{G}_a$ 而是 $\OmX$，于是我们得到一个投影
\[ p_\alpha\colon U^{\Th}\to\OmX. \]
把所有 $p_\alpha$ 相加，我们得到投影 $U^{\Th}\to\OmX$。过渡到分类空间并从 $X$ 取映射叠，这给出映射
\[ \Bun_U^{\Th}(\Fq)\to H^1(X,\OmX), \]
而 Serre 对偶给出同构 $H^1(X,\OmX)=\Fq$。最后，Whittaker 特征是复合
\[ \Bun_U^{\Th}(\Fq)\to\Fq\xrightarrow{\tr}\Fp\hookrightarrow\C^\times, \]
其中用到 $\C$ 中单位根的标准选取。
\end{remark}

如上所述，该定理是数十年来几何 Langlands 纲领研究的成果。几何 Langlands 纲领的起点是 Drinfeld 的观察：在某些情形，可以取一个 $\Gd$-局部系统，并在另一侧产生相应的自守形式，参看 Drinfeld（1983）。

\srcpage{10}
为此，Drinfeld 实际上构造了 $\Bun_G$ 上的一个 $\ell$ 进层 $\cF$，然后通过 Grothendieck 的层-函数字典得到 $\Bun_G(\Fq)$ 上的函数 $\tr_\cF$：对任意点 $x\in\Bun_G(\Fq)$，对应 $\tr(\Frob_x|\cF_{\bar x})$，即（$x$ 上方一个几何点 $\bar x$ 处的）Frobenius 迹。这提示，在几何上，$X_{\overline{\F}_q}$ 上的 $\Gd$-局部系统与 $\Bun_{G,\overline{\F}_q}$ 上的可构造层之间可能存在关系。

非常粗略地说，定理 1.7 的证明现在如下：
\begin{enumerate}
\item 把几何 Langlands 对应表述为范畴等价，大致说它把 $\Bun_{G,\overline{\F}_q}$ 上的 étale $\ell$ 进层
\[ D_{\et}(\Bun_{G,\overline{\F}_q},\overline{\Q}_\ell) \]
与 $\Gd$-局部系统的叠 $\LocSys^{\mathrm{geom}}_{\Gd}$ 上的拟凝聚层
\[ D_{\qc}(\LocSys^{\mathrm{geom}}_{\Gd}) \]
联系起来，其中 $\LocSys^{\mathrm{geom}}_{\Gd}$ 是 $X_{\overline{\F}_q}$ 上 $\Gd$-局部系统的叠。
\item 证明在这两个范畴上施行「取 Frobenius 迹」的操作，就得到所需的向量空间同构。
\item 对其他域上的光滑射影曲线、以及其他层论也表述几何 Langlands 对应，并证明所得的猜想本质上不依赖于所选定的设置。
\item 对特征 $0$ 的域上的曲线、在 D-模的层论中证明几何 Langlands 猜想。
\end{enumerate}
在下一节，我们将给出几何 Langlands 等价的一般性介绍；在随后的各节中，将聚焦于其表述与证明的不同方面。

\begin{remark}
在最初的数域情形，这一策略在第一步就无可挽回地失效，因为不存在 $X_{\overline{\F}_q}$ 的类比。长期以来，人们认为几何 Langlands 纲领所研究的结构不适用于最初的数域情形。然而，Fargues（2025）提议在一个特殊的「曲线」——Fargues–Fontaine 曲线上研究几何 Langlands 猜想，它在某种意义上把 $p$ 进数 $\Q_p$ 几何化。结果表明，这给出 $p$ 进域上局部 Langlands 猜想的一个精确版本，参看 Fargues–Scholze（2024）。在实数上也有平行的故事，参看 Ben-Zvi–Nadler（2013）、Scholze（2024）。局部算术 Langlands 目前的许多研究都建立在几何 Langlands 的范式之上。然而，在整体数域情形仍然需要真正新的想法。我们注意到，所需同构似乎极不可能通过「在某个范畴等价上取 Frobenius 迹」的类比而得到\footnote{在数域与 3 维流形的类比中，这最终归结为：有限域上的函数域对应以圆周为底的 3 维流形纤维化，而数域只对应 3 维流形，没有这种纤维化。}。
\end{remark}

\srcpage{11}
\textbf{致谢}。本报告只不过触及了皮毛，我为其中众多的遗漏、不精确与未竟之处致歉。关于几何 Langlands 对应之证明的另一篇生动得多的综述，我推荐 Ben-Zvi（2026）。我很高兴感谢许多数学家，他们教给我大量关于 Langlands 纲领、特别是几何 Langlands 纲领的知识，他们的洞见与远见塑造了我的理解：Johannes Anschütz、David Ben-Zvi、Arthur-César le Bras、Ana Caraiani、Dustin Clausen、Vladimir Drinfeld、Matt Emerton、Laurent Fargues、Dennis Gaitsgory、Toby Gee、Linus Hamann、David Hansen、Eugen Hellmann、Tim Kuppel、Vincent Lafforgue、Gérard Laumon、Lucas Mann、David Nadler、Michael Rapoport、Sam Raskin、Timo Richarz、Juan Esteban Rodríguez Camargo、Yiannis Sakellaridis、Germán Stefanich、Akshay Venkatesh、Kari Vilonen、Cong Xue、Zhiwei Yun、Xinwen Zhu，以及许多其他人。特别感谢 Dennis Gaitsgory 与 Sam Raskin 在准备本报告期间回答了我的许多问题；对于任何可能的误述，我谨致歉意。
